\begin{onehalfspace}

Questo capitolo introduce l'algoritmo che dà a questa tesi il proprio oggetto. Si apre con i richiami di calcolo quantistico indispensabili a seguirlo --- il qubit, la sovrapposizione, le porte unitarie e il meccanismo di interferenza --- di cui l'algoritmo costituisce l'applicazione più celebre.

L'\textbf{algoritmo di Shor}, proposto da Peter Shor nel 1994 \cite{shor1994}, è un algoritmo quantistico che risolve il problema della \textbf{fattorizzazione intera} di un numero $N$ in tempo \textbf{polinomiale} nel numero di bit di $N$:
\begin{equation}
    T_{\text{Shor}}(N) = O\!\left((\log N)^3\right)
    \label{eq:shor_complexity}
\end{equation}

Questo risultato costituisce uno \textbf{speedup esponenziale} rispetto al miglior algoritmo classico noto -- il \ac{GNFS} -- che richiede tempo sub-esponenziale (si veda l'equazione~\eqref{eq:gnfs_complexity}), del tutto impraticabile per i numeri di grandi dimensioni usati in crittografia. È considerato il risultato algoritmico più dirompente dell'informatica quantistica, poiché la fattorizzazione è il fondamento matematico del sistema crittografico RSA, che protegge la quasi totalità delle comunicazioni digitali moderne.

L'algoritmo si articola in tre fasi: un \textit{preprocessing} classico che riconduce il problema a trovare il periodo di una funzione modulare; un nucleo quantistico basato sulla \ac{QFT} e sulla \ac{QPE}, che trova quel periodo in modo esponenzialmente più efficiente di qualunque strategia classica; e una \textit{post-elaborazione} classica che estrae i fattori di $N$ tramite il calcolo del massimo comune divisore. Dopo i richiami si introduce il problema che l'algoritmo risolve e la sua rilevanza crittografica, per poi svilupparne nel dettaglio la struttura.

\section{Richiami di Calcolo Quantistico}
\label{chap:fondamenti}

Questa sezione richiama i concetti di meccanica quantistica e di calcolo quantistico necessari a seguire il resto della tesi, riportando per esteso quelli su cui i capitoli successivi si appoggiano operativamente --- la sfera di Bloch, le matrici di Pauli e la loro algebra. La trattazione completa, con i quattro postulati, l'insieme esteso delle porte e il modello di complessità, è nella Sezione~\ref{app:fondamenti}.

\subsection{Qubit, Sovrapposizione e Sfera di Bloch}

L'unità fondamentale del calcolo quantistico è il \textbf{qubit}, descritto da un vettore unitario nello spazio di Hilbert $\mathcal{H} = \mathbb{C}^2$. In notazione di Dirac, lo stato generico è la \textbf{sovrapposizione}
\begin{equation}
    \ket{\psi} = \alpha\ket{0} + \beta\ket{1}, \qquad \alpha, \beta \in \mathbb{C}, \quad |\alpha|^2 + |\beta|^2 = 1,
\end{equation}
dove $|\alpha|^2$ e $|\beta|^2$ sono le probabilità di misurare rispettivamente $0$ e $1$ (\textbf{regola di Born}).

Riparametrizzando le due ampiezze complesse tramite angoli reali --- il che è lecito perché la normalizzazione toglie un grado di libertà e la fase globale è inosservabile --- ogni stato puro si scrive
\begin{equation}
    \ket{\psi} = \cos\!\left(\frac{\theta}{2}\right)\ket{0} + e^{i\phi}\sin\!\left(\frac{\theta}{2}\right)\ket{1},
    \qquad \theta \in [0,\pi],\ \phi \in [0,2\pi),
    \label{eq:bloch_param}
\end{equation}
e si visualizza come un punto sulla superficie della \textbf{sfera di Bloch}.\footnote{La sfera deve il nome al fisico svizzero-americano Felix Bloch (1905--1983), Premio Nobel per la Fisica nel 1952 per lo sviluppo della spettroscopia a risonanza magnetica nucleare (\ac{NMR}). La stessa notazione $T_1$ e $T_2$ per i tempi di decoerenza proviene dal formalismo NMR di Bloch.} I poli sono $\ket{0}$ e $\ket{1}$, l'equatore le sovrapposizioni equiprobabili con fasi diverse.

\begin{figure}[H]
    \centering
    \includegraphics[width=0.42\textwidth]{bloch_sphere.png}
    \caption[La sfera di Bloch]{La sfera di Bloch: ogni stato puro di un qubit corrisponde a un punto sulla superficie della sfera unitaria, parametrizzato dagli angoli $\theta$ e $\varphi$ della~\eqref{eq:bloch_param}. I poli nord e sud rappresentano gli stati base $\ket{0}$ e $\ket{1}$; i punti sull'equatore corrispondono a sovrapposizioni equiprobabili. È la rappresentazione a cui si farà riferimento nel Capitolo~\ref{chap:rumore} per descrivere l'azione dei canali di rumore, che deformano la sfera anziché spostare punti su di essa.}
    \label{fig:bloch}
\end{figure}

\subsection{Sistemi Multi-Qubit ed Entanglement}

Un registro di $n$ qubit vive nello spazio $\mathbb{C}^{2^n}$, ottenuto per \textbf{prodotto tensoriale}. La dimensione esponenziale rende possibile rappresentare ampiezze e correlazioni non disponibili a un registro classico di $n$ bit, ma non implica da sola un vantaggio computazionale: alcune famiglie di stati e circuiti quantistici ammettono infatti descrizioni classiche efficienti. Due qubit sono \textbf{entangled} quando il loro stato congiunto non è fattorizzabile nel prodotto degli stati individuali; l'esempio canonico è lo stato di Bell $\ket{\Phi^+} = \tfrac{1}{\sqrt{2}}(\ket{00} + \ket{11})$, per il quale misure nella stessa base producono esiti perfettamente correlati, senza consentire comunicazione istantanea.

\subsubsection{L'Entanglement come Risorsa Fisica Fragile}

L'entanglement è una risorsa importante per Shor: durante l'esponenziazione modulare il registro di controllo si correla con quello di lavoro, permettendo alla successiva Quantum Fourier Transform di rendere osservabile la periodicità. Un qubit reale, però, non è perfettamente isolato: l'accoppiamento con l'ambiente riduce le coerenze dello stato --- la \textbf{decoerenza}~\cite{zurek2003}. Il tempo caratteristico della perdita di fase è $T_2$ (\textit{dephasing time}); nei dispositivi superconduttori è spesso dello stesso ordine di poche centinaia di durate di una porta a due qubit. Questo confronto segnala un budget di profondità severo, ma non determina una soglia universale nella dimensione di $N$: compilazione, parallelismo, tempi di idle, qualità dei gate e architettura fault-tolerant concorrono tutti al limite discusso nel Capitolo~\ref{chap:rumore}.

\subsection{Porte, Misura e Interferenza}

Le operazioni sui qubit sono \textbf{matrici unitarie} ($U^\dagger U = I$), quindi reversibili. La Tabella~\ref{tab:porte_essenziali} raccoglie quelle su cui il resto della tesi si appoggia; l'insieme completo, con le derivazioni, è nella Sezione~\ref{app:fondamenti}.

\begin{table}[H]
\centering
\begin{tabular}{llc l}
\toprule
\textbf{Porta} & \textbf{Matrice} & \textbf{Qubit} & \textbf{Effetto} \\
\midrule
$X$ & $\begin{pmatrix} 0 & 1 \\ 1 & 0 \end{pmatrix}$ & 1 & \textit{bit flip}: scambia $\ket{0}$ e $\ket{1}$ \\[10pt]
$Y$ & $\begin{pmatrix} 0 & -i \\ i & 0 \end{pmatrix}$ & 1 & bit flip e phase flip insieme ($\propto XZ$) \\[10pt]
$Z$ & $\begin{pmatrix} 1 & 0 \\ 0 & -1 \end{pmatrix}$ & 1 & \textit{phase flip}: $\ket{1} \mapsto -\ket{1}$ \\[10pt]
$H$ & $\tfrac{1}{\sqrt{2}}\begin{pmatrix} 1 & 1 \\ 1 & -1 \end{pmatrix}$ & 1 & genera la sovrapposizione uniforme \\[10pt]
$R_\phi$ & $\begin{pmatrix} 1 & 0 \\ 0 & e^{i\phi} \end{pmatrix}$ & 1 & rotazione di fase; $T = R_{\pi/4}$, $S = R_{\pi/2}$ \\[10pt]
CNOT & $\begin{pmatrix} 1&0&0&0 \\ 0&1&0&0 \\ 0&0&0&1 \\ 0&0&1&0 \end{pmatrix}$ & 2 & nega il \textit{target} se il \textit{control} è $\ket{1}$ \\
\bottomrule
\end{tabular}
\caption[Le porte quantistiche essenziali]{Le porte su cui poggia il resto dell'elaborato. Le tre matrici di Pauli descrivono i canali di rumore del Capitolo~\ref{chap:rumore} e gli errori che i codici correttori dei Capitoli~\ref{chap:qec}--\ref{chap:surface} devono diagnosticare; $H$ e $R_\phi$ costruiscono la \ac{QFT}; la \ac{CNOT} è la porta a due qubit dominante nel conteggio delle risorse, ed è quella a cui è associato l'errore più severo.}
\label{tab:porte_essenziali}
\end{table}

L'identità e le tre matrici di Pauli formano una base dello spazio degli operatori lineari $2\times2$: ogni operatore d'errore su un singolo qubit si può quindi decomporre in tale base. Poiché $Y=iXZ$, un codice che soddisfa le condizioni di correzione per gli errori di Pauli singoli corregge, per linearità, anche le loro combinazioni --- proprietà su cui si fonda l'intera parte \ac{QEC}. Valgono inoltre le relazioni
\begin{equation}
\begin{aligned}
    X^2 = Y^2 = Z^2 &= I, \\
    XY = iZ, \quad YZ = iX, \quad ZX &= iY, \\
    \{X,Y\} = \{Y,Z\} = \{X,Z\} &= 0,
\end{aligned}
    \label{eq:algebra_pauli}
\end{equation}
dove $\{A,B\} = AB+BA$ è l'anticommutatore. Nella teoria degli stabilizzatori, un errore è rilevato quando anticommuta con almeno uno dei generatori misurati; è il vettore completo di questi esiti di commutazione a costituire la sindrome. Nei codici \ac{CSS}, le famiglie di stabilizzatori $X$ e $Z$ permettono di diagnosticare separatamente le due componenti di Pauli, come discusso nel Capitolo~\ref{chap:qec}.

Insieme alle porte a singolo qubit, la \ac{CNOT} forma un insieme universale: ogni operazione unitaria su $n$ qubit si approssima a piacere componendo elementi di questo insieme. La \textbf{misura} in base computazionale collassa lo stato $\ket{\psi}=\sum_x\alpha_x\ket{x}$ sull'esito $x$ con probabilità $|\alpha_x|^2$: è un'operazione non unitaria e irreversibile. Il meccanismo che rende utile tutto ciò è l'\textbf{interferenza}: un algoritmo efficiente è costruito perché le ampiezze degli esiti sbagliati si cancellino (interferenza distruttiva) e quelle degli esiti corretti si sommino (costruttiva). È esattamente ciò che la QFT realizza nell'algoritmo di Shor (Capitolo~\ref{chap:shor}) per far emergere il periodo.

\end{onehalfspace}
